\chapter{中期进展报告}


\section{已完成的任务及取得的成果}

按照开题阶段制定的研究计划，本项目目前已按进度完成以下任务并取得相应成果：
\begin{itemize}
    \item 完成交通时空预测领域的调研工作，系统梳理了代表性预测方法、公开基准数据集与评价指标；
    \item 完成基线方法的实验结果复现，并在此基础上开展对比实验；
    \item 完成研究方法主体框架的搭建，并在基准公开数据集上完成实验，获得初步实验结果；
    \item 完成参数敏感性实验和模型参数的可视化分析，初步验证了所提方法的有效性；
    \item 在前期的多元时间序列预测研究中，已有一篇论文被 ACM CIKM 2025 国际顶级会议录用发表。
\end{itemize}

\section{尚需完成的任务}

在前期工作基础上，后续尚需完成以下任务：
\begin{itemize}
    \item 进一步完善方法设计，重点优化多尺度空间构造、专家路由与静动态特征融合策略，提升模型在复杂路网和不同交通状态下的预测能力；
    \item 扩展实验规模，补充更全面的对比实验、消融实验和可视化分析，进一步验证方法的有效性；
    \item 将交通时空预测研究成果整理成论文并完成投稿，同时按计划完成学位论文的撰写。
\end{itemize}

\section{存在的问题}

当前研究主要存在以下问题：
\begin{itemize}
    \item 所提方法在部分数据集或评价指标上与同期先进方法相比仍存在一定差距，模型结构和训练策略仍有优化空间；
    \item 稀疏门控混合专家在训练稳定性和专家负载均衡方面仍有改进余地；
    \item 模型对节假日、突发事件等非规律性交通场景的建模能力有待进一步加强。
\end{itemize}

\section{拟采取的办法}

针对上述问题，拟采取以下办法：
\begin{itemize}
    \item 针对部分数据集或指标上的性能差距，细化多尺度空间图构造和静动态双专家的路由与融合策略，并通过消融实验定位关键模块，进行有针对性的改进；
    \item 引入专家负载均衡机制和合适的正则化策略，提升稀疏激活训练过程的稳定性；
    \item 补充更多数据集和场景下的对比实验与可视化分析，结合典型案例分析改进模型在非规律性场景下的泛化能力；
    \item 尽快完成实验结果整理与论文撰写，按计划推进论文投稿和学位论文写作。
\end{itemize}
